% ============================================================
% CHAPTER 1: INTRODUCTION
% ============================================================
\customchapter{Introduction}

This chapter establishes the context, motivation, and objectives behind the development of RivetDB. We examine the limitations of existing in-memory data stores, formulate the problem statement, and outline the project scope and software development methodology.

\section{Background and Motivation}
In the modern digital era, the volume, velocity, and variety of data generated by applications have reached unprecedented levels. Systems across domains such as financial technology, high-frequency trading, e-commerce, real-time analytics, and the Internet of Things (IoT) require data to be ingested, processed, and queried with latencies measured in microseconds. Traditional disk-based relational database management systems (RDBMS) such as PostgreSQL and MySQL, while offering robust ACID (Atomicity, Consistency, Isolation, Durability) guarantees, are fundamentally constrained by the physical limitations of block storage devices. Even with the advent of Non-Volatile Memory Express (NVMe) solid-state drives, the latency overhead of disk I/O operations, page faulting, and buffer pool management introduces unacceptable delays for ultra-low-latency use cases.

This paradigm shift birthed the widespread adoption of In-Memory Database Systems (IMDBs). By storing the entire operational dataset within the system's main memory (RAM), IMDBs bypass disk access paths entirely, relying on direct memory access and CPU cache hierarchies. This architectural choice yields orders-of-magnitude improvements in read and write throughput. The ecosystem of in-memory key-value stores has historically been bifurcated into two dominant architectural models, represented by Redis and Memcached.

\textbf{Redis (Remote Dictionary Server)}: Introduced in 2009, Redis revolutionized caching and message brokering by offering a rich set of data structures (Lists, Sets, Sorted Sets, Hashes) rather than simple strings. However, its core command execution model is strictly single-threaded. This design choice, made by its creator to avoid the complexities of thread synchronization, deadlocks, and race conditions, means that a single Redis instance can only ever utilize a single CPU core. In modern cloud environments where multi-core processors are ubiquitous, this represents a severe utilization bottleneck.

\textbf{Memcached}: Predating Redis, Memcached was designed explicitly for multi-threaded, high-throughput caching of simple string key-value pairs. It leverages a pool of worker threads to process concurrent requests, allowing it to scale linearly with available CPU cores. However, Memcached offers almost no advanced data structures and lacks any persistence mechanism, limiting its use cases strictly to volatile caching.

Furthermore, both Redis and Memcached are implemented in the C programming language. Decades of software engineering history have demonstrated that manual memory management in C is extraordinarily prone to critical security vulnerabilities, such as buffer overflows, use-after-free errors, and null-pointer dereferences. The Rust programming language emerged to solve exactly this problem. By enforcing a strict ownership model and borrow checker at compile-time, Rust guarantees memory safety and freedom from data races without the runtime overhead of a garbage collector.

\section{Problem Statement}
The prevailing in-memory data store ecosystem forces software architects to accept significant compromises, creating a void in high-performance infrastructure tooling. Developers must constantly choose between the rich features and safety of a single-threaded system (Redis) or the multi-core scalability of a limited, volatile caching layer (Memcached). 

Furthermore, standard deployments of these systems lack native support for modern access patterns. Features such as SQL querying, infinite multi-tenancy isolation, and robust time-series data aggregation are typically bolted on as afterthoughts, forcing reliance on complex, unstable, or commercially licensed third-party modules. This fragmentation leads to increased operational complexity and fragility in production environments.

Finally, the reliance on C and C++ for legacy IMDBs presents an ever-present vector for memory-safety vulnerabilities in critical data infrastructure. In an era where data breaches cost organizations millions, building core data stores on unsafe memory foundations is an unacceptable risk. There is a pressing need for a system that combines the multi-core scalability of Memcached, the rich data structures of Redis, and the guaranteed memory safety of modern systems programming languages like Rust.

\section{Objective}
This project aims to design, implement, and benchmark \textbf{RivetDB}, a concurrent, memory-safe, and highly scalable in-memory database built entirely in Rust. The primary objectives are:

\textbf{1. Concurrent Architecture Implementation:} To implement a multi-threaded command execution engine using fine-grained locking mechanisms (DashMap) and an asynchronous network I/O runtime (Tokio) that scales predictably and linearly across multiple CPU cores, effectively solving the single-thread bottleneck of legacy systems.

\textbf{2. Absolute Memory Safety Guarantees:} To leverage the Rust programming language's strict compile-time borrow checker to categorically eliminate memory corruption vulnerabilities, segmentation faults, and data races, ensuring enterprise-grade stability without a garbage collector.

\textbf{3. Rich Data Model Integration:} To engineer and support a wide array of native data structures, far beyond simple key-value strings. This includes advanced structures such as Lists, Sets, Sorted Sets, Hashes, JSON documents, Time-Series data, and Bloom Filters, all natively integrated into the command parser.

\textbf{4. Embedded SQL Query Engine:} To design and embed a native SQL-like parsing and execution engine allowing complex declarative queries over the key-value store, bridging the gap between NoSQL performance and relational querying flexibility.

\textbf{5. Strict Multi-Tenancy Isolation:} To design an infinite namespace isolation mechanism to allow multiple applications or tenants to share a single database instance with strict memory boundary enforcement and zero cross-tenant data leakage.

\textbf{6. Non-Blocking Persistence:} To implement a low-overhead, fully asynchronous Append-Only File (AOF) persistence layer to ensure durability without blocking the main event loop, achieving the durability of disk databases with the speed of RAM.

\textbf{7. Operational Observability:} To embed a production-grade Prometheus-compatible HTTP metrics server exposing real-time telemetry including per-command throughput, latency distributions, memory utilization, key-type breakdowns, hot-key detection, and a slow-query log for operational debugging.

\textbf{8. Authentication and Access Control:} To implement a connection-level authentication framework using the standard \texttt{AUTH} command protocol, with configurable password enforcement and per-connection session state tracking.

\textbf{9. TOML-Based Runtime Configuration:} To design a hierarchical configuration system using TOML files and command-line argument overrides, allowing operators to tune server bindings, memory limits, eviction policies, persistence behavior, and security parameters without recompilation.

\textbf{10. Schema Registry and Type-Safe Keys:} To engineer a JSON Schema validation subsystem allowing administrators to define strict type contracts for key namespaces, enforcing data integrity at the database tier via \texttt{TSET}, \texttt{TGET}, \texttt{TVALIDATE}, and \texttt{TKEYS} commands.

\section{Scope of the Project}
The scope of RivetDB is bounded to a single-node, in-memory execution environment. It is designed to maximize vertical scaling on modern multi-core servers. The project includes the full implementation of the TCP network layer, RESP protocol parsing, the core concurrent storage engine, eight native data types, asynchronous disk persistence, a Prometheus metrics HTTP server, connection authentication, and a TOML-based configuration system. 

Out of scope for this phase is distributed multi-node clustering (horizontal scaling) and consensus algorithms like Raft or Paxos. Additionally, while the system persists data to an AOF file for recovery, it is strictly an in-memory-first database; datasets that exceed available RAM are managed via eviction policies rather than paging to disk.

\section{Software Development Life Cycle (SDLC)}
The development of RivetDB followed an Agile iterative Software Development Life Cycle (SDLC). Unlike traditional Waterfall models, an iterative approach was essential due to the highly exploratory nature of concurrent systems programming in Rust.

The phases included continuous Requirement Analysis, where the bottlenecks of existing systems were evaluated. This was followed by System Design, focusing on lock-free and fine-grained locking data structures. The Implementation phase involved heavy use of Rust's Tokio ecosystem. Crucially, the Testing \& Benchmarking phase was deeply integrated into the iterations; performance regressions were identified immediately, allowing the design to adapt dynamically.

\begin{figure}[htbp]
    \centering
    \begin{tikzpicture}[
        node distance=1.5cm and 2cm,
        box/.style={draw, rectangle, rounded corners, minimum width=3cm, minimum height=1cm, align=center},
        arrow/.style={-stealth, thick},
        dashed_arrow/.style={-stealth, dashed, thick}
    ]
    
    % Define the blocks
    \node[box] (A) {Requirement\\Analysis};
    \node[box, right=of A] (B) {System\\Design};
    \node[box, below=of B] (C) {Implementation\\in Rust};
    \node[box, left=of C] (D) {Testing \&\\Benchmarking};
    \node[box, left=of D] (E) {Deployment \&\\Review};
    
    % Draw the connecting arrows
    \draw[arrow] (A) -- (B);
    \draw[arrow] (B) -- (C);
    \draw[arrow] (C) -- (D);
    \draw[arrow] (D) -- (E);
    \draw[dashed_arrow] (E) |- node[near start, left] {Iterative Feedback} (A);
    
    \end{tikzpicture}
    \caption{Agile Software Development Life Cycle for RivetDB}
    \label{fig:sdlc_diagram}
\end{figure}

\section{System Requirements}
To successfully compile, deploy, and benchmark the RivetDB system, the following hardware and software prerequisites must be met.

\subsection{Hardware Requirements}
\begin{itemize}[leftmargin=2.5em]
    \item \textbf{Processor}: Multi-core CPU (Minimum 4 cores, 8+ recommended for optimal concurrent scaling).
    \item \textbf{Memory (RAM)}: Minimum 4 GB (8 GB or more recommended depending on the target dataset size).
    \item \textbf{Storage}: Solid State Drive (SSD) recommended for optimal AOF persistence write throughput.
\end{itemize}

\subsection{Software Requirements}
\begin{itemize}[leftmargin=2.5em]
    \item \textbf{Operating System}: Linux (Ubuntu 20.04/22.04 recommended), macOS, or Windows (via WSL2).
    \item \textbf{Compiler Toolchain}: Rust (\texttt{rustc} version 1.70.0 or higher) and the \texttt{cargo} package manager.
    \item \textbf{Client Tools}: standard \texttt{redis-cli} or any Redis-compatible client library (e.g., \texttt{redis-py}, \texttt{ioredis}).
\end{itemize}

\section{Organization of Project Report}
The remainder of this report is organized as follows:
\begin{itemize}[leftmargin=2.5em]
    \item \textbf{Chapter 2: Literature Survey} provides a comprehensive review of 10 foundational research papers detailing existing in-memory systems, the memory wall problem, and concurrency paradigms.
    \item \textbf{Chapter 3: System Design} details the high-level architecture, flowcharts, block diagrams, network layer design, and the integration of the SQL engine.
    \item \textbf{Chapter 4: Implementation} provides an in-depth code-level analysis of the RivetDB source code.
    \item \textbf{Chapter 5: Performance Optimization} discusses the specific profiling tools used and the tuning strategies applied to maximize throughput.
    \item \textbf{Chapter 6: Results and Discussion} presents the findings of our extensive benchmark testing, applications, and advantages.
    \item \textbf{Chapter 7: Problems and Drawbacks} analyzes the architectural limitations and missing features of the current implementation.
    \item \textbf{Chapter 8: Conclusion and Future Scope} summarizes the accomplishments of the project and outlines the roadmap for future enhancements.
\end{itemize}

